\documentclass[11pt,a4paper]{article}
\usepackage{amsmath,amssymb,amsthm,bm}
\usepackage{geometry}
\geometry{margin=1in}
\usepackage{hyperref}
\usepackage{graphicx}
\usepackage{setspace}
\usepackage{booktabs}
\onehalfspacing

\title{The Parasitic Threshold and Boundary Logic:\\ Defense Pressure Before Modern Immunity}
\author{Steve Bico Mujjabi, MD\thanks{msbico@gmail.com}\\
VuraLabs, Kampala, Uganda\\
Independent researcher\\
ORCID: \href{https://orcid.org/0009-0001-0556-5516}{0009-0001-0556-5516}}
\date{August 2026}

\begin{document}
\maketitle

\begin{abstract}
Parasitism becomes meaningful only after replication, retention, and container inheritance exist. This manuscript therefore moves the parasitic threshold after the protocell integration gate (Paper~9). In CDFD terms, a parasitic sequence or pathway is an internal flux drain that consumes activated resources while failing to maintain the host network or its boundary. The key claim is conditional and testable: once copying and compartment persistence are coupled, exploitative variants create selection pressure for specificity, exclusion, proofreading, degradation, and boundary control. The paper does not claim literal modern viruses, immune systems, or oncology at the dawn of life --- it identifies the earliest functional problem that later defense systems solve.
\end{abstract}

\noindent\textbf{Keywords:} parasitic threshold; boundary logic; prebiotic defense; replication; origins of life.

\section*{Preprint note}
This is a computational research article. The companion script produces numerical diagnostics reported in the Computational Results section. It is not a narrative review and does not claim historical proof of a unique abiogenesis pathway.

\section*{Symbol Conventions}

\noindent See Paper~1, \S Symbol Conventions, for the base CDFD and Greek-symbol tables. This paper introduces $\Phi_{\mathrm{host}}$, $\Phi_{\mathrm{loss}}$, and the host-survival inequality (\S4).

\section{Introduction}
Prior protocell-integration work established the minimum machinery required for protocellular evolution: information must copy with bounded error, and container success must covary with internal chemistry. Once that machinery exists, exploitation becomes unavoidable in principle. Any variant that uses the replication or catalytic apparatus without contributing to persistence can grow faster than cooperative chemistry over short horizons.

This is the parasitic threshold. It is not a statement that modern viruses existed before cells. It is a statement about a failure mode of coupled autocatalytic and replicative systems, adjacent to classic error-threshold and major-transition logic~\cite{Eigen1971,MaynardSmith1997}.

\section{Parasitic Flux Drain}
Let the cooperative autocatalytic/replicative network have operational flux $\Phi_{\mathrm{host}}$. A parasitic component consumes resources at rate $\Phi_{\mathrm{loss}}$ while contributing little or nothing to retention, repair, catalytic closure, or boundary persistence. The effective host flux becomes
\begin{equation}
 \Phi_{\mathrm{eff}}=\Phi_{\mathrm{host}}-\Phi_{\mathrm{loss}}.
\end{equation}
The adaptive ratio becomes
\begin{equation}
 \Psi_s=
 \left(\frac{\Phi_{\mathrm{eff}}}{C_{\mathrm{mem}}}\right)S M_s.
\end{equation}
If $\Phi_{\mathrm{loss}}$ grows faster than host repair or exclusion, the system crosses from productive replication to starvation, washout, or loss of inherited organization.

\section{Boundary Logic}
The response cannot be only ``more flux.'' Extra substrate can amplify the parasite as easily as the host. A surviving protocell-like system needs constraint changes:
\begin{itemize}
 \item selective permeability at the boundary,
 \item catalytic specificity inside the compartment,
 \item degradation or sequestration of noncontributing templates,
 \item copying rules that preserve useful heredity while limiting runaway freeloaders.
\end{itemize}
These are early functional analogues of defense logic. They are not yet modern immune systems.

\section{Mathematical Gate}
A simple host-survival inequality is
\begin{equation}
\Gamma_{\mathrm{host}}Q_{\mathrm{host}}R_{\mathrm{ret}}
>
\Gamma_{\mathrm{parasite}}Q_{\mathrm{parasite}}R_{\mathrm{access}},
\end{equation}
where $R_{\mathrm{ret}}$ measures host retention in the container and $R_{\mathrm{access}}$ measures parasite access to host resources. Boundary evolution acts by lowering $R_{\mathrm{access}}$ or increasing the cost of parasitic replication.

The remaining symbols are operational. $\Gamma_{\mathrm{host}}$ is the productive host replication or catalytic-throughput rate, $Q_{\mathrm{host}}$ is host copying fidelity or functional retention per copy, $\Gamma_{\mathrm{parasite}}$ is the parasite replication or drain rate, and $Q_{\mathrm{parasite}}$ is parasite persistence per copy. The inequality is not a calibrated law; it is a gate that says cooperative persistence must outpace exploitative access after boundary effects are included.

The falsifiable implication is direct: in a model or experiment with heritable compartments, increasing noncontributing replicators should collapse cooperative persistence unless some specificity or exclusion mechanism is introduced.

\section{Validation Design}
The experimental target is a compartmental copying system in which host-like templates or catalysts contribute to persistence while parasite-like variants consume the same activated resources without contributing to boundary survival. The controls should hold total template concentration, activated substrate, compartment size distribution, and environmental cycling fixed while varying only parasite fraction and boundary specificity.

Useful observables include host-template fraction after serial transfer, compartment survival, product yield per compartment, leakage, and the covariance between internal chemistry and boundary growth or division. The CDFD expectation is a threshold response: increasing parasite access should sharply reduce host persistence unless selective permeability, degradation, sequestration, or copying specificity is present. Failure to see such a threshold under otherwise heritable conditions would weaken this paper's claim.

\section{Computational Results: Spatial Boundary Gate}
The final release expands the companion script beyond a well-mixed
host--parasite ODE. A second diagnostic places host-like and parasite-like
fields on a two-dimensional compartment grid. The boundary gate modifies
parasite access to the interior while the top of the domain receives resource
inflow. This remains a toy stress test, but it makes the physical point
visible: boundary specificity is not just a scalar penalty; it changes where
resource-consuming variants can spatially persist.

The two-dimensional panel should be read as a falsification target. If parasite
suppression appears only in a well-mixed equation but disappears once diffusion,
leakage, and boundary access are represented spatially, the boundary-logic
claim is too weak. Conversely, if gated compartments retain host-dominant
interiors across parameter sweeps while open boundaries do not, the model earns
a stronger experimental target.

\section{Conclusion}
The parasitic threshold is the first crisis created by successful protocell integration. Once copying, retention, and compartment inheritance exist, exploitative variants force the system to evolve boundary logic. Photochemical expansion is treated as a later energy-input channel (Paper~11) rather than a prerequisite for the parasitic-threshold argument.

\section{Reproducibility Note}
The companion script \texttt{supplementary\_ool\_10.py} tests a minimal
host--parasite competition gate with explicit retention, boundary specificity,
and a two-dimensional spatial boundary diagnostic. It writes
\path{parasitic_threshold_boundary_logic.png} under \path{outputs/paper_10/},
plus well-mixed and spatial CSV/JSON summaries. No notebook is required.

\begin{figure}[!htbp]
\centering
\includegraphics[width=\linewidth,height=0.30\textheight,keepaspectratio]{outputs/paper_10/parasitic_threshold_boundary_logic.png}
\caption{Parasitic-threshold diagnostic. The toy model compares host persistence with and without boundary specificity when a noncontributing replicator accesses the same resource stream, then shows a two-dimensional boundary-gate stress test. The figure is a mechanism diagnostic, not a claim about modern viruses.}
\end{figure}

\section*{Related literature}
This manuscript is self-contained within the Part~II origins-of-life sequence. Shared CDFD/CDFL symbols ($\Phi$, $C$, $S$, $M_s$, $\Psi_s$) are used as bookkeeping only; they do not require Part~I physics results. Field anchors are cited in the body and bibliography.

\section*{Competing Interests}
The author declares no competing interests.

\section*{Funding}
This research received no external funding.

\section*{Data and Code Availability}
Companion scripts, generated figures, and numerical outputs for this manuscript are included in the public CDFD Part II release archive (\url{https://github.com/bampita-bico/CDFD-Part-II-Release}; Zenodo DOI \texttt{10.5281/zenodo.20264779}). See the Reproducibility Note above for the paper-local script path.

\section*{Author Contributions}
Conceptualization, formal analysis, computational diagnostics, and manuscript preparation: Steve Bico Mujjabi

\bibliographystyle{plain}
\bibliography{ool_references}
\end{document}
